\documentclass[12pt]{article}
\usepackage[utf8]{inputenc}
\usepackage[T1]{fontenc}
\usepackage{hyperref}
\title{Tonjon Language Resources for K-12 Curriculum Development}
\author{Compiled by Claude Sonnet 5 (claude-sonnet-5) for Andr\'as Kornai}
\date{}
\begin{document}
\maketitle

\section{Introduction}
Tonjon (ISO 639-3 tjn, Glottocode tonj1246) is an extinct Western Mande language of C\^ote d'Ivoire, in the Jogo-Jeri branch. It was reportedly spoken by blacksmiths in the Djimini Senoufo community and is closely related to Ligbi. Wikipedia gives its extinction as since the 1950s. I found no independent documentation of the last speakers. In practice, no teaching or documentation resources turned up in my search.

\section{Primary Educational Resources}
None found. There are no primers, textbooks or school materials for Tonjon.

\section{Linguistic and Reference Materials}
\begin{itemize}
\item Tonjon language (Wikipedia): \url{https://en.wikipedia.org/wiki/Tonjon_language}
\item Glottolog entry: \url{https://glottolog.org/resource/languoid/id/tonj1246.bigmap.html}
\item Ethnologue subgroup page: \url{https://www.ethnologue.com/subgroup/974/}
\item Related language, Djimini Senoufo (still spoken): \url{https://en.wikipedia.org/wiki/Djimini_language}
\item Jogo languages: \url{https://en.wikipedia.org/wiki/Jogo_languages}
\end{itemize}
These pages are overview and classification material only. I did not find a grammar, wordlist or dictionary.

\section{Religious and Community Resources}
None found.

\section{Digital and Audio Resources}
None found.

\section{Curriculum Development Context}
A native-language curriculum is not feasible from what I found. The most that can be done is a heritage-awareness unit on the history of the Djimini blacksmith community and on the Mande languages of C\^ote d'Ivoire. Before building anything, one would need to find archival field notes through linguists who work on Jogo-Jeri and Ligbi. I did not verify that any such notes exist.
\end{document}